\documentclass[UTF8,zihao=5,a4paper]{ctexart}

% 老师反馈（2026-07-15，留在源文件中供后续推进时核对）：
% 1. 文献调研需要提取本课题的关键创新点，重点区分已有的单跳/多跳风险传播
%    与我们拟研究的远距离作用；“超距”只能在有交通机理和统计检验支撑后使用。
% 2. 需要明确研究对象。当前先收敛为混合交通中的碰撞/冲突风险，避免把事故、
%    拥堵和一般运行风险混在同一个定义中。

\usepackage[top=1.85cm,bottom=1.85cm,left=2.1cm,right=2.1cm]{geometry}
\usepackage{amsmath}
\usepackage{booktabs}
\usepackage{enumitem}
\usepackage[hidelinks]{hyperref}
\usepackage{float}
\usepackage{tikz}
\usetikzlibrary{arrows.meta,positioning}

\setlength{\parindent}{2em}
\setlength{\parskip}{0.12em}
\linespread{1.08}
\renewcommand{\arraystretch}{1.15}
\setlist[enumerate]{leftmargin=2.2em,itemsep=1pt,topsep=2pt,parsep=0pt}
\ctexset{
  section={format=\large\bfseries,beforeskip=10pt,afterskip=5pt},
  subsection={format=\normalsize\bfseries,beforeskip=7pt,afterskip=3pt}
}
\setlength{\textfloatsep}{7pt}
\setlength{\intextsep}{6pt}
\setlength{\abovecaptionskip}{4pt}
\setlength{\belowcaptionskip}{2pt}

\title{\vspace{-1.2cm}\textbf{混合自动驾驶风险传播研究\\阶段工作回顾与本周进展}}
\author{禹尧珅}
\date{7.9--7.14}

\begin{document}

\pagestyle{plain}
\maketitle
\vspace{-1.8em}

\section{前期工作回顾}

\subsection{前期文献调研与可行性判断}

前期我一共整理了22篇文献，主要分为风险事件与安全指标、混合自动驾驶传播与控制、闭环仿真与场景测试、人类行为交互与轨迹预测四类。我当时主要想确认三件事：现有数据能不能识别局部风险，风险传播能不能用图来表示，以及后面能不能用闭环仿真检验干预效果。

\textbf{混合交通交互建模。}Wang等~\cite{wang2024safety}将机理模型与高斯过程结合，用实车实验数据学习人工驾驶车辆的速度变化和预测不确定性，再把模型放进预测控制器中。文中还用稀疏高斯过程和动态预测降低在线计算量。\textbf{这篇文章对我最直接的帮助是：}分析风险传播时要保留自动驾驶车辆与人工驾驶车辆之间的交互，后面做干预时也要同时看安全变化和控制代价。

Gong和Zhu~\cite{gong2024hinf}建立了包含不同跟驰方式和反应延迟的车辆跟驰模型，用车辆自身状态和前车加速度描述运动，并推导串稳定条件，再通过线性矩阵不等式设计鲁棒控制和$H_{\infty}$控制。数值结果表明，自动驾驶车辆占比和反应延迟都会改变扰动抑制能力。\textbf{这篇文章提醒我：}正式实验不能只跑一个车辆占比，还要更换随机种子，比较传播距离、延迟和影响范围是否发生变化。

\textbf{风险图建模。}Wang等~\cite{wang2024conflict}提出冲突风险图，把路口中容易发生冲突的位置设为节点，并按空间拓扑连接这些节点；随后构建基于Transformer的图卷积网络，用图卷积提取局部空间特征，用Transformer处理较长时间范围的依赖。\textbf{我准备借用其中的图建模思路：}把文章中的“风险位置节点”改成“车辆风险事件节点”，再用车辆关系和时间顺序连接前后事件。

\textbf{闭环仿真验证。}Fu等~\cite{fu2025trafficmcts}提出基于分组蒙特卡洛树搜索的闭环交通流生成框架，并用社会价值取向描述不同车辆的合作倾向和驾驶风格。车辆会根据每一步的环境状态重新决策，形成连续的多车交互过程。\textbf{这一点正好对应后面的反事实实验：}保持初始条件不变，只修改某个行为并重新运行，再比较前后两次仿真的风险链。

\textbf{综合这些文献，我认为这个方向是可以继续做的。}数据方面，车辆轨迹中已经有速度、加速度、间距和前后车关系；方法方面，可以参考现有风险图来设计事件节点和传播边；验证方面，闭环仿真可以在固定初始条件下做干预对照。

\subsection{前期实验基础}

前期我已经在DRIFT/Flow环境中跑过混合交通合流场景。仿真可以设置自动驾驶车辆占比，并输出每辆车的时间、位置、速度、加速度、车头时距和前车关系等数据。我也完成了规则型风险检测，并跑通了一次合流场景。

当时的分析主要统计风险次数和全局指标，同一个持续风险会在连续多帧中被重复记录，事件之间也没有连接起来。因此，这周我一边补充文献，一边尝试把逐帧检测整理成事件片段，再构建传播边。

\section{本周工作（7.9--7.14）}

\subsection{文献调研}

这周我在前期调研的基础上又补了两类文献：\textbf{因果风险建模}和\textbf{传播过程量化}。主要目的是弄清楚当前传播边应该怎么解释，以及后面还要统计哪些指标。

\textbf{因果风险建模。}Lyu等~\cite{lyu2025causal}给出了“因果发现、数据重构、实时预测”的完整流程：先用结构不可知模型从交通变量中寻找因果关系，再根据因果图和时空关系重构输入，最后用图注意力网络预测快速路分流区的冲突风险。\textbf{我从这篇文章里确定了一个边界：}现在根据时间、距离和车辆关系得到的边只能算候选关联，关键边还要放回闭环仿真中重新验证。

\textbf{传播过程量化。}Guan等~\cite{guan2025jam}先从微观跟驰角度分析冲击波，再用随机游走描述混合交通中拥堵波的传播，并推导受影响车辆数量等指标；随后设计缓入策略，并用数值仿真检查抑制效果。\textbf{对我来说最有用的是指标设计：}除了事件数量，还要统计传播延迟、传播距离、影响车辆数和传播深度，干预时可以优先选择下游影响范围较大的事件。

看完这两类工作后，后面的思路基本确定为：\textbf{先用可解释规则构建候选传播图，再把关键关系放回闭环仿真中做对照。}

\subsection{小实验验证}

\textbf{（1）数据与风险事件。}这次小实验用的是一次Flow合流场景产生的车辆轨迹，自动驾驶车辆占比为20\%，仿真时长为25秒。我按加速度、碰撞时间和车头时距，从每一帧车辆状态中识别急刹、严重急刹、低碰撞时间、临界碰撞时间、低车头时距和近碰事件。其中，车辆接近前车时的碰撞时间为
\begin{equation}
T_i=\frac{h_i}{-\Delta v_i}, \qquad \Delta v_i<0,
\label{eq:ttc}
\end{equation}
式中，$h_i$为车辆$i$与前车的间距，$\Delta v_i$为相对速度。当前阈值包括：急刹加速度不高于$-3\,\mathrm{m/s^2}$、严重急刹加速度不高于$-6\,\mathrm{m/s^2}$、低碰撞时间不超过2秒、低车头时距不超过1秒。

为了让不同类型的事件可以放在一起比较，代码把每条检测的风险分数$q_i$归一化到$[0,1]$。对于碰撞时间或车头时距等“数值越小越危险”的指标，分数定义为
\begin{equation}
q_z=\min\left\{1,\max\left(0,\frac{\theta-z}{\theta}+0.5\right)\right\},
\label{eq:risk-inverse}
\end{equation}
其中，$z$是当前指标值，$\theta$是对应的风险阈值。对于急刹事件，分数按减速度相对预设上限的比例计算；急刹和严重急刹分别使用$8\,\mathrm{m/s^2}$和$10\,\mathrm{m/s^2}$作为归一化上限。分数越接近1，表示这条检测在当前规则下越严重。这个分数只用于比较相对严重程度，不是真实碰撞概率。

急刹类事件的具体计算式为
\begin{equation}
q_a=\min\left\{1,\max\left(\frac{|a|}{a_{\mathrm{cap}}},\frac{a_{\mathrm{th}}}{a_{\mathrm{cap}}}\right)\right\},
\label{eq:risk-braking}
\end{equation}
其中，$a_{\mathrm{th}}$为触发阈值的绝对值，$a_{\mathrm{cap}}$为归一化上限。近碰事件同时满足低碰撞时间和急刹条件，其风险分数取碰撞时间分数与0.8中的较大值。

\textbf{（2）事件片段合并。}逐帧检测会重复记录持续风险。设第$i$个检测为$e_i=(r_i,v_i,c_i,t_i,q_i)$，分别表示运行编号、车辆编号、风险类型、发生时间和风险分数。相邻检测满足
\begin{equation}
r_i=r_j,\quad v_i=v_j,\quad c_i=c_j,\quad 0<t_j-t_i\leq 0.4\ \mathrm{s}
\label{eq:merge}
\end{equation}
时，将其合并为同一风险片段。片段的风险分数取内部检测的最大值，并记录起止时间、持续时间和原始检测数。

\textbf{（3）传播边筛选。}对时间上先后出现且满足同一车辆、前后车或同一区域关系的两个风险片段建立候选边。时间差限制为5秒，局部空间距离限制为80米。候选边$i\rightarrow j$的传播分数按现有代码计算为
\begin{align}
S_{ij}&=\min\left\{1,\max\left(0,w_t w_d w_r w_c\right)\right\}, \label{eq:score}\\
w_t&=\max\left(0.1,1-\frac{\Delta t_{ij}}{5}\right),\quad
w_d=\max\left(0.1,1-\frac{\min(d_{ij},80)}{80}\right), \notag\\
w_r&=\max(q_i,q_j). \notag
\end{align}
其中，$\Delta t_{ij}$为时间差，$d_{ij}$为空间距离，$q_i$和$q_j$为两端事件的风险分数；$w_c$为关系权重，同车、前后车、同一区域和远距离同车道分别取0.9、1.0、0.7和0.5。为控制图的规模，每个源事件在每类关系中只保留得分最高的后续事件。车道字段缺失时，不保留远距离同车道边。

传播分数$S_{ij}$也在$[0,1]$内。两个事件间隔越短、距离越近、风险越高，而且车辆关系越直接，得分就越高。这个分数只用来排列候选边和筛选直接边，不能当作传播概率或因果强度。

\textbf{（4）节点角色与干预优先级。}设事件$i$的入边强度、出边强度和下游风险增量分别为
\begin{equation}
I_i=\sum_{j\rightarrow i}S_{ji},\qquad
O_i=\sum_{i\rightarrow j}S_{ij},\qquad
G_i=\operatorname{mean}_{i\rightarrow j}\max(q_j-q_i,0).
\label{eq:role-base}
\end{equation}
令$\mathcal{N}(x_i)=x_i/\max_k x_k$表示同一次运行内的最大值归一化，则三个节点角色分数为
\begin{align}
R_i^{\mathrm{source}}&=\mathcal{N}(O_i), \notag\\
R_i^{\mathrm{amplifier}}&=\mathcal{N}\!\left(\min(I_i,O_i)(1+G_i)\right), \label{eq:roles}\\
R_i^{\mathrm{absorber}}&=\mathcal{N}\!\left(\frac{I_i}{1+O_i}\right). \notag
\end{align}
这三个分数主要是为了帮助我读图。风险源分数高，表示该事件后面连接了较强的传播边；放大器分数高，表示它接收了上游影响，又继续向下游传播，而且下游风险可能增加；吸收器分数高，表示它接收到的影响较强，但后续传播较弱。这些角色只针对当前这一次运行，不代表某辆车长期固定属于某一类。

在直接图上向下遍历最多三层后，干预优先级定义为
\begin{equation}
P_i=R_i^{\mathrm{source}}(1+n_i)\left(1+\sum_{e\in E_i}S_e\right),
\label{eq:intervention}
\end{equation}
其中，$n_i$是从事件$i$能够到达的下游事件数，$E_i$是遍历过程中经过的直接传播边。$P_i$越高，说明该事件连接的下游范围越大、累计传播分数越高，后面可以优先拿它做反事实实验。这个优先级只用于挑选实验对象，干预有没有效果仍要重新运行闭环仿真。

\begin{figure}[H]
\centering
\begin{tikzpicture}[
  node distance=5mm,
  box/.style={draw=black,rectangle,align=center,minimum width=2.18cm,minimum height=0.72cm,inner sep=3pt},
  flow/.style={-{Latex[length=2mm]},thick}
]
\node[box] (data) {闭环仿真\\车辆轨迹};
\node[box,right=of data] (event) {逐帧风险\\事件识别};
\node[box,right=of event] (episode) {连续检测\\合并为片段};
\node[box,right=of episode] (edge) {候选边构建\\直接边筛选};
\node[box,right=of edge] (role) {节点角色\\干预候选};
\draw[flow] (data) -- (event);
\draw[flow] (event) -- (episode);
\draw[flow] (episode) -- (edge);
\draw[flow] (edge) -- (role);
\end{tikzpicture}
\caption{小实验的风险传播分析流程}
\label{fig:flow}
\end{figure}

\textbf{（5）当前结果。}这次实验一共得到309条逐帧风险检测。按照式\eqref{eq:merge}合并后形成16个风险片段，随后生成37条候选传播边，并从中筛选出18条直接传播边。具体结果见表\ref{tab:results}。

\begin{table}[H]
\centering
\caption{真实Flow合流数据的分析结果}
\label{tab:results}
\begin{tabular}{lr}
\toprule
指标 & 结果 \\
\midrule
逐帧风险检测数 & 309 \\
风险事件片段数 & 16 \\
候选传播边数 & 37 \\
直接传播边数 & 18 \\
平均传播延迟 & 1.467秒 \\
平均传播距离 & 13.679米 \\
平均传播分数 & 0.507 \\
\bottomrule
\end{tabular}
\end{table}

309条记录来自连续仿真帧，所以不能把它理解成309个独立事件。合并后的16个片段更接近完整的风险过程。37条候选边最后保留18条，保留率为48.6\%，图中的重复联系和较间接的联系明显减少。18条直接边中有9条前后车关系、5条同区域关系和4条同车关系。就这次样本来看，比较清楚的联系主要出现在前后车之间。

平均传播延迟为1.467秒，平均传播距离为13.679米，说明目前筛出的关系主要发生在较短时间和局部范围内。平均传播分数0.507是18条直接边分数的算术平均，只能在同一套评分规则下比较，目前还没有设置高、中、低等级。由于现在只有一个场景、20\%自动驾驶车辆占比和一次运行，车道字段也没有完整记录，所以这组结果主要用来检查分析流程。跨车道传播、一般规律和干预效果还需要后续实验支持。

\section{下周计划（7.15--7.21）}

\begin{enumerate}
\item \textbf{优化实验设计。}我准备先补全车辆类型、车道和前后车关系等字段，再确定不同自动驾驶车辆占比和随机种子的对照设置。
\item \textbf{完善传播分析。}继续检查事件定义和传播边筛选方法，补充传播深度、影响车辆数等指标，并选出适合做反事实实验的事件。
\item \textbf{推进文档与论文工作。}把研究问题、方法、实验设置和阶段结果整理清楚，同时列出各部分还缺少的证据。
\end{enumerate}

\begingroup
\footnotesize
\begin{thebibliography}{99}
\bibitem{wang2024conflict}
T. Wang, Y.-E. Ge, Y. Wang, C. G. Prato, W. Chen, and Y. Niu,
``A conflict risk graph approach to modeling spatio-temporal dynamics of intersection safety,''
\emph{Transportation Research Part C: Emerging Technologies}, 2024.

\bibitem{lyu2025causal}
N. Lyu, T. Xie, J. Wen, and Y. Wang,
``A real-time conflict risk prediction modeling based on causal inference and graph-model: Applied to multi-configuration expressway diversion zones,''
\emph{Transportation Research Part C: Emerging Technologies}, 2025.

\bibitem{guan2025jam}
H. Guan, X. Chen, and Q. Meng,
``Jam propagation in mixed traffic of autonomous and human-driven vehicles: A random walk-based analysis,''
\emph{Transportation Research Part C: Emerging Technologies}, vol. 180, p. 105310, 2025.

\bibitem{wang2024safety}
J. Wang, Y. V. Pant, L. Zhao, M. Antkiewicz, and K. Czarnecki,
``Enhancing safety in mixed traffic: Learning-based modeling and efficient control of autonomous and human-driven vehicles,''
\emph{IEEE Transactions on Intelligent Transportation Systems}, vol. 25, no. 9, 2024.

\bibitem{fu2025trafficmcts}
Z. Fu, L. Wen, P. Cai, D. Fu, S. Mao, and B. Shi,
``TrafficMCTS: A closed-loop traffic flow generation framework with group-based Monte Carlo tree search,''
\emph{IEEE Transactions on Intelligent Transportation Systems}, vol. 26, no. 10, 2025.

\bibitem{gong2024hinf}
Y. Gong and W.-X. Zhu,
``Modeling and robust H-infinity control synthesis of the CAV-HDV heterogeneous traffic system with different car-following modes,''
\emph{IEEE Transactions on Intelligent Transportation Systems}, vol. 25, no. 10, 2024.
\end{thebibliography}
\endgroup

\end{document}
